\section{Threats to Validity}

\textit{Construct validity.} Our ground truth is a line diff: two findings are the same when they share a rule and the old line maps to the new one, unchanged or edited with a similarity of at least 0.6. A finding whose line is rewritten past that threshold, or moved out of a deleted block, counts as removed plus introduced. We checked 100 findings drawn in five strata from the CodeGuard replays. One author labelled the 60 persisting ones; an LLM rater labelled all 100 without seeing the ground truth or the author's labels; the author then reviewed the 14 findings where the LLM rater said the ground truth was wrong, and agreed on all of them. Re-running the rule showed that three of these, complexity and nesting findings on an unchanged \texttt{def} line, were new after all: the commit had pushed the function over its limit. With the other 11 corrected, the ground truth held for 59 of 60 persisting findings (98.3\%, Wilson 95\% CI 91.1--99.7) and 30 of 40 introduced ones (75.0\%, 59.8--85.8). Every error on an introduced finding was one some scheme had matched to an old finding, so the hidden rates in Table~\ref{tab:rq1} are upper bounds. SonarQube's falls from 1.47\% to an estimated 0.24\% (0.07--0.66), though each corrected rate rests on only 6 to 12 sampled findings. The 27 introduced findings the LLM rater confirmed as new had no human check, and Bandit's own findings were not sampled.

\textit{Internal validity.} We measure ports of each tool's identity method, not the tools. The SonarQube port matched the unmodified Java tracker on 20,000 generated cases, and the Semgrep port reproduced all 716 IDs Semgrep emitted in two replays; the Bandit and GitHub ports agree with \texttt{bandit -b} and with GitHub's \texttt{fingerprints.ts}. SonarQube's branch-mode move detection is modelled only for files git reports as renamed. Some CodeGuard messages contain line numbers, which penalizes methods that compare messages, so we also report Bandit's method on Bandit's own findings, without B101 (91\% of them, nearly all in tests).

\textit{External validity.} The corpus is 35 open-source web applications in Python and JavaScript or TypeScript, each replayed over 200 consecutive commit pairs. Consecutive pairs are the easy case; a real baseline is compared with code many commits later, which we do not measure. A few repositories can dominate a pooled rate, so we report intervals over repositories and per-repository medians. In RQ2 each transformation is applied alone, to one finding, while real commits combine several.
